% GENERATED FILE. Edit canonical lessons, then run npm run generate:books.
% canonical-source-hash: fnv1a64:43ebc28550298919
% canonical-lessons: UR-C342-vot, UR-C342-intikhab, UR-C342-jamhuriyat, UR-C342-ummidvar, UR-C342-thopna

\chapter{Citizens and Votes}
\label{ch:ur-citizens-and-votes}
\hlchaptermodality{342}

\begin{chapteropening}
\textbf{By the end of this chapter:} \emph{I can talk about voting and the rights of citizens.}

Complete the last lesson of chapter 342: I can talk about voting and the rights of citizens.
\end{chapteropening}

\section[\texorpdfstring{\ur{ووٹ} (voṭ)}{ووٹ (voṭ)}]{\ur{ووٹ} (voṭ) — a vote}

\label{lesson:UR-C342-vot}



\begin{quote}
Before the new one: say the Urdu for the power to rule, then the Urdu for to hold to be.
\end{quote}

\subsection*{You'll want to know: \ur{ووٹ}}
\textbf{\ur{ووٹ}} — \emph{voṭ} — \textquotedblleft{}a vote\textquotedblright{}.

\begin{quote}
\ur{میں} \ur{نے} \ur{صبح} \ur{اپنا} \ur{ووٹ} \ur{دیا۔}
\end{quote}

\begin{quote}
\emph{maiṅ ne subah apnā voṭ diyā.}
\end{quote}

\begin{quote}
I cast my vote in the morning.
\end{quote}

Masculine.

\begin{cousinweb}
\textbf{\ur{ووٹ}} (\emph{voṭ}) is borrowed from English \emph{vote}.
\end{cousinweb}

\begin{grammarlens}[title={What you've built}]
One more word for giving reasons and opinions.
\end{grammarlens}

\subsection*{Your turn}
\begin{itemize}
\raggedright
  \item \emph{Say it:} \emph{voṭ}
  \item \emph{Say it:} \emph{voṭ}, once more
  \item \emph{Say it:} \emph{maiṅ ne subah apnā voṭ diyā.} — I cast my vote in the morning.
  \item \emph{Recall:} say the Urdu for the power to rule, then the Urdu for to hold to be, then say \emph{voṭ} again
\end{itemize}

\begin{tcolorbox}[breakable,colback=teal!4,colframe=teal!35!black,title={Before you move on}]
What does \ur{ووٹ} mean? (\textquotedblleft{}A vote\textquotedblright{}.) Say it once more.
\end{tcolorbox}
\section[\texorpdfstring{\ur{انتخاب} (intikhāb)}{انتخاب (intikhāb)}]{\ur{انتخاب} (intikhāb) — an election; a choice}

\label{lesson:UR-C342-intikhab}



\begin{quote}
Before the new one: say the Urdu for to hold to be, then the Urdu for a vote.
\end{quote}

\subsection*{You'll want to know: \ur{انتخاب}}
\textbf{\ur{انتخاب}} — \emph{intikhāb} — \textquotedblleft{}an election; a choice\textquotedblright{}.

\begin{quote}
\ur{اگلے} \ur{سال} \ur{انتخاب} \ur{ہو} \ur{گا۔}
\end{quote}

\begin{quote}
\emph{agle sāl intikhāb ho gā.}
\end{quote}

\begin{quote}
There will be an election next year.
\end{quote}

Masculine.

\begin{cousinweb}
\textbf{\ur{انتخاب}} (\emph{intikhāb}) is usually traced, through Persian, to Arabic \emph{intikhāb}, \textquotedblleft{}choosing\textquotedblright{}.
\end{cousinweb}

\begin{grammarlens}[title={What you've built}]
One more word for giving reasons and opinions.
\end{grammarlens}

\subsection*{Your turn}
\begin{itemize}
\raggedright
  \item \emph{Say it:} \emph{intikhāb}
  \item \emph{Say it:} \emph{intikhāb}, once more
  \item \emph{Say it:} \emph{agle sāl intikhāb ho gā.} — There will be an election next year.
  \item \emph{Recall:} say the Urdu for to hold to be, then the Urdu for a vote, then say \emph{intikhāb} again
\end{itemize}

\begin{tcolorbox}[breakable,colback=teal!4,colframe=teal!35!black,title={Before you move on}]
What does \ur{انتخاب} mean? (\textquotedblleft{}An election; a choice\textquotedblright{}.) Say it once more.
\end{tcolorbox}
\section[\texorpdfstring{\ur{جمہوریت} (jamhūriyat)}{جمہوریت (jamhūriyat)}]{\ur{جمہوریت} (jamhūriyat) — democracy}

\label{lesson:UR-C342-jamhuriyat}



\begin{quote}
Before the new one: say the Urdu for a vote, then the Urdu for an election.
\end{quote}

\subsection*{You'll want to know: \ur{جمہوریت}}
\textbf{\ur{جمہوریت}} — \emph{jamhūriyat} — \textquotedblleft{}democracy\textquotedblright{}.

\begin{quote}
\ur{جمہوریت} \ur{میں} \ur{ہر} \ur{شہری} \ur{کی} \ur{رائے} \ur{ضروری} \ur{ہے۔}
\end{quote}

\begin{quote}
\emph{jamhūriyat meṅ har shahrī kī rāe zarūrī hai.}
\end{quote}

\begin{quote}
In a democracy every citizen's opinion matters.
\end{quote}

Feminine.

\begin{cousinweb}
\textbf{\ur{جمہوریت}} (\emph{jamhūriyat}) is usually traced, through Persian, to Arabic \emph{jumhūriyya}, \textquotedblleft{}a republic\textquotedblright{}, built on \emph{jumhūr}, \textquotedblleft{}the people at large\textquotedblright{}.
\end{cousinweb}

\begin{grammarlens}[title={What you've built}]
One more word for giving reasons and opinions.
\end{grammarlens}

\subsection*{Your turn}
\begin{itemize}
\raggedright
  \item \emph{Say it:} \emph{jamhūriyat}
  \item \emph{Say it:} \emph{jamhūriyat}, once more
  \item \emph{Say it:} \emph{jamhūriyat meṅ har shahrī kī rāe zarūrī hai.} — In a democracy every citizen's opinion matters.
  \item \emph{Recall:} say the Urdu for a vote, then the Urdu for an election, then say \emph{jamhūriyat} again
\end{itemize}

\begin{tcolorbox}[breakable,colback=teal!4,colframe=teal!35!black,title={Before you move on}]
What does \ur{جمہوریت} mean? (\textquotedblleft{}Democracy\textquotedblright{}.) Say it once more.
\end{tcolorbox}
\section[\texorpdfstring{\ur{امیدوار} (ummīdvār)}{امیدوار (ummīdvār)}]{\ur{امیدوار} (ummīdvār) — a candidate (for a post or a seat)}

\label{lesson:UR-C342-ummidvar}



\begin{quote}
Before the new one: say the Urdu for an election, then the Urdu for democracy.
\end{quote}

\subsection*{You'll want to know: \ur{امیدوار}}
\textbf{\ur{امیدوار}} — \emph{ummīdvār} — \textquotedblleft{}a candidate (for a post or a seat)\textquotedblright{}.

\begin{quote}
\ur{اس} \ur{انتخاب} \ur{میں} \ur{تین} \ur{امیدوار} \ur{ہیں۔}
\end{quote}

\begin{quote}
\emph{is intikhāb meṅ tīn ummīdvār haiṅ.}
\end{quote}

\begin{quote}
There are three candidates in this election.
\end{quote}

Masculine or feminine, by the person meant.

\begin{cousinweb}
\textbf{\ur{امیدوار}} (\emph{ummīdvār}) is usually traced to Persian \emph{ummīdvār}, \textquotedblleft{}hopeful\textquotedblright{}, built on \textbf{\ur{امید}}, \textquotedblleft{}hope\textquotedblright{}, with \emph{-vār}, \textquotedblleft{}having\textquotedblright{}: the one who hopes to be chosen.
\end{cousinweb}

\begin{grammarlens}[title={What you've built}]
One more word for giving reasons and opinions.
\end{grammarlens}

\subsection*{Your turn}
\begin{itemize}
\raggedright
  \item \emph{Say it:} \emph{ummīdvār}
  \item \emph{Say it:} \emph{ummīdvār}, once more
  \item \emph{Say it:} \emph{is intikhāb meṅ tīn ummīdvār haiṅ.} — There are three candidates in this election.
  \item \emph{Recall:} say the Urdu for an election, then the Urdu for democracy, then say \emph{ummīdvār} again
\end{itemize}

\begin{tcolorbox}[breakable,colback=teal!4,colframe=teal!35!black,title={Before you move on}]
What does \ur{امیدوار} mean? (\textquotedblleft{}A candidate (for a post or a seat)\textquotedblright{}.) Say it once more.
\end{tcolorbox}
\section[\texorpdfstring{\ur{تھوپنا} (thopnā)}{تھوپنا (thopnā)}]{\ur{تھوپنا} (thopnā) — to impose; to force (on someone)}

\label{lesson:UR-C342-thopna}



\begin{quote}
Before the new one: say the Urdu for democracy, then the Urdu for a candidate.
\end{quote}

\subsection*{You'll want to know: \ur{تھوپنا}}
\textbf{\ur{تھوپنا}} — \emph{thopnā} — \textquotedblleft{}to impose; to force (on someone)\textquotedblright{}.

\begin{quote}
\ur{حکومت} \ur{نے} \ur{لوگوں} \ur{پر} \ur{نیا} \ur{ٹیکس} \ur{تھوپ} \ur{دیا۔}
\end{quote}

\begin{quote}
\emph{hukūmat ne logoṅ par nayā ṭaiks thop diyā.}
\end{quote}

\begin{quote}
The government forced a new tax on people.
\end{quote}

Past, with \ur{نے} (he or she …): \ur{اس} \ur{نے} \ur{اپنی} \ur{رائے} \ur{تھوپی}, \emph{us ne apnī rāe thopī}.

\begin{cousinweb}
\textbf{\ur{تھوپنا}} (\emph{thopnā}) is inherited: CDIAL lists the Hindi cousin \emph{thopnā}, \textquotedblleft{}to cover, to heap\textquotedblright{}, under \emph{*stōpyate}, a form it reconstructs (CDIAL 13723); the sense \textquotedblleft{}to force on someone\textquotedblright{} is a later one our sources do not trace.
\end{cousinweb}

\begin{grammarlens}[title={What you've built}]
One more word for giving reasons and opinions.
\end{grammarlens}

\subsection*{Your turn}
\begin{itemize}
\raggedright
  \item \emph{Say it:} \emph{thopnā}
  \item \emph{Say it:} \emph{thopnā}, once more
  \item \emph{Say it:} \emph{hukūmat ne logoṅ par nayā ṭaiks thop diyā.} — The government forced a new tax on people.
  \item \emph{Recall:} say the Urdu for democracy, then the Urdu for a candidate, then say \emph{thopnā} again
\end{itemize}

\begin{tcolorbox}[breakable,colback=teal!4,colframe=teal!35!black,title={Before you move on}]
What does \ur{تھوپنا} mean? (\textquotedblleft{}To impose; to force (on someone)\textquotedblright{}.) Say it once more.
\end{tcolorbox}
